% Seminar 2 - Pagina 3. Fragment pentru % Seminar 2 - Pagina 3. Fragment pentru % Seminar 2 - Pagina 3. Fragment pentru % Seminar 2 - Pagina 3. Fragment pentru \input{seminars/pages/seminar2-pagina3-ro}.
% Fara preambul si fara intrerupere fortata de pagina.
% NOTE EDITORIALE:
% - Fotografia este numerotata cu 3 in coltul din dreapta jos, desi
%   cererea mentioneaza pagina 2. Acest fragment este deci pagina 3.
% - In manuscris, f(X) si f^{-1}(Y) sunt definite la inceputul exercitiului;
%   anumite semne din dreptul cuantificatorilor sunt suprascrise.
% - Exemplul din dreapta despre f:Z->{0,1} contine intervale/notatii
%   greu lizibile. Sunt redate numai citirile provizorii [-1,0] si [1,2].
%   Manuscrisul noteaza rezultatele ca 1 si 0, nu ca multimi singleton.
% - Sub punctul b) exista etape taiate; au fost omise. In sens strict,
%   si demonstratia de la a) foloseste nejustificat acelasi antecedent x
%   pentru egalitatea imaginilor, desi f nu este presupusa injectiva.
% - Egalitatea la b) pentru o familie vida ar necesita in plus f(A)=B;
%   manuscrisul nu discuta aceasta exceptie.
%   afirmatia f(x) in f(intersectie) nu implica x in toate X_i fara
%   injectivitate: e nevoie de un martor potrivit in intersectie.
% - Pagina se termina cu inceputul punctului c) din demonstratie,
%   fara alte randuri; nu a fost inventata continuarea.

\subsection*{Seminar 2 - Pagina 3}

\noindent\fbox{4}\quad Fie $f:A\longrightarrow B$, $X\subseteq A$, $Y\subseteq B$.
\[
 f(X)=\{f(x)\mid x\in X\}, \qquad
 f^{-1}(Y)=\{a\in A\mid f(a)\in Y\}.
\]
\noindent (Preimaginea lui $Y$; fibra lui $Y$.)

\medskip
\noindent Fie $f:A\longrightarrow B$,
$(X_i)_{i\in I}\subseteq\mathcal{P}(A)$,
$(Y_j)_{j\in J}\subseteq\mathcal{P}(B)$.
\begin{enumerate}
\item[a)]
\[
 f\left(\bigcup_{i\in I}X_i\right)
  =\bigcup_{i\in I}f(X_i).
\]
\item[b)]
\[
 f\left(\bigcap_{i\in I}X_i\right)
  \subseteq\bigcap_{i\in I}f(X_i).
\]
\noindent Dacă $f$ este injectivă, avem egalitate.
\item[c)]
\[
 f^{-1}\left(\bigcup_{j\in J}Y_j\right)
  =\bigcup_{j\in J}f^{-1}(Y_j).
\]
\item[d)]
\[
 f^{-1}\left(\bigcap_{j\in J}Y_j\right)
  =\bigcap_{j\in J}f^{-1}(Y_j).
\]
\end{enumerate}

% Exemplul din marginea dreapta a manuscrisului (citire provizorie).
\noindent\underline{Ex.} $f:\mathbb Z\longrightarrow\{0,1\}$,
\[
 f(x)=\begin{cases}
  0, & x>0,\\
  1, & x\leq 0.
 \end{cases}
 \qquad
 f([-1,0])=1,\quad f([1,2])=0.
\]

\medskip
\noindent\underline{Sol.}

\noindent $a)$ Fie
$ f(x)\in f\left(\bigcup_{i\in I}X_i\right)$.
Atunci $x\in\bigcup_{i\in I}X_i$, deci există $i_1\in I$
astfel încât $x\in X_{i_1}$, de unde
\[
 f(x)\in f(X_{i_1})\subseteq\bigcup_{i\in I}f(X_i).
\]
\noindent Rezultă incluziunea $\subseteq$.

\smallskip
\noindent Reciproc, dacă
$f(x)\in\bigcup_{i\in I}f(X_i)$, există $i_1\in I$
a.î.\ $f(x)\in f(X_{i_1})$. Prin urmare,
\[
 x\in X_{i_1}\subseteq\bigcup_{i\in I}X_i
 \quad\Longrightarrow\quad
 f(x)\in f\left(\bigcup_{i\in I}X_i\right).
\]

\medskip
\noindent $b)$ Dacă
$f(x)\in f\left(\bigcap_{i\in I}X_i\right)$, atunci
$x\in X_i$ pentru orice $i\in I$, de unde
\[
 f(x)\in f(X_i)\quad(\forall i\in I)
 \quad\Longrightarrow\quad
 f(x)\in\bigcap_{i\in I}f(X_i).
\]
\noindent Dacă $f$ este injectivă, din
$f(x)\in\bigcap_{i\in I}f(X_i)$ rezultă
$f(x)\in f(X_i)$ pentru orice $i$, deci
$x\in\bigcap_{i\in I}X_i$.

\medskip
\noindent $c)$
% Demonstratia punctului c) continua pe urmatoarea pagina de manuscris.
.
% Fara preambul si fara intrerupere fortata de pagina.
% NOTE EDITORIALE:
% - Fotografia este numerotata cu 3 in coltul din dreapta jos, desi
%   cererea mentioneaza pagina 2. Acest fragment este deci pagina 3.
% - In manuscris, f(X) si f^{-1}(Y) sunt definite la inceputul exercitiului;
%   anumite semne din dreptul cuantificatorilor sunt suprascrise.
% - Exemplul din dreapta despre f:Z->{0,1} contine intervale/notatii
%   greu lizibile. Sunt redate numai citirile provizorii [-1,0] si [1,2].
%   Manuscrisul noteaza rezultatele ca 1 si 0, nu ca multimi singleton.
% - Sub punctul b) exista etape taiate; au fost omise. In sens strict,
%   si demonstratia de la a) foloseste nejustificat acelasi antecedent x
%   pentru egalitatea imaginilor, desi f nu este presupusa injectiva.
% - Egalitatea la b) pentru o familie vida ar necesita in plus f(A)=B;
%   manuscrisul nu discuta aceasta exceptie.
%   afirmatia f(x) in f(intersectie) nu implica x in toate X_i fara
%   injectivitate: e nevoie de un martor potrivit in intersectie.
% - Pagina se termina cu inceputul punctului c) din demonstratie,
%   fara alte randuri; nu a fost inventata continuarea.

\subsection*{Seminar 2 - Pagina 3}

\noindent\fbox{4}\quad Fie $f:A\longrightarrow B$, $X\subseteq A$, $Y\subseteq B$.
\[
 f(X)=\{f(x)\mid x\in X\}, \qquad
 f^{-1}(Y)=\{a\in A\mid f(a)\in Y\}.
\]
\noindent (Preimaginea lui $Y$; fibra lui $Y$.)

\medskip
\noindent Fie $f:A\longrightarrow B$,
$(X_i)_{i\in I}\subseteq\mathcal{P}(A)$,
$(Y_j)_{j\in J}\subseteq\mathcal{P}(B)$.
\begin{enumerate}
\item[a)]
\[
 f\left(\bigcup_{i\in I}X_i\right)
  =\bigcup_{i\in I}f(X_i).
\]
\item[b)]
\[
 f\left(\bigcap_{i\in I}X_i\right)
  \subseteq\bigcap_{i\in I}f(X_i).
\]
\noindent Dacă $f$ este injectivă, avem egalitate.
\item[c)]
\[
 f^{-1}\left(\bigcup_{j\in J}Y_j\right)
  =\bigcup_{j\in J}f^{-1}(Y_j).
\]
\item[d)]
\[
 f^{-1}\left(\bigcap_{j\in J}Y_j\right)
  =\bigcap_{j\in J}f^{-1}(Y_j).
\]
\end{enumerate}

% Exemplul din marginea dreapta a manuscrisului (citire provizorie).
\noindent\underline{Ex.} $f:\mathbb Z\longrightarrow\{0,1\}$,
\[
 f(x)=\begin{cases}
  0, & x>0,\\
  1, & x\leq 0.
 \end{cases}
 \qquad
 f([-1,0])=1,\quad f([1,2])=0.
\]

\medskip
\noindent\underline{Sol.}

\noindent $a)$ Fie
$ f(x)\in f\left(\bigcup_{i\in I}X_i\right)$.
Atunci $x\in\bigcup_{i\in I}X_i$, deci există $i_1\in I$
astfel încât $x\in X_{i_1}$, de unde
\[
 f(x)\in f(X_{i_1})\subseteq\bigcup_{i\in I}f(X_i).
\]
\noindent Rezultă incluziunea $\subseteq$.

\smallskip
\noindent Reciproc, dacă
$f(x)\in\bigcup_{i\in I}f(X_i)$, există $i_1\in I$
a.î.\ $f(x)\in f(X_{i_1})$. Prin urmare,
\[
 x\in X_{i_1}\subseteq\bigcup_{i\in I}X_i
 \quad\Longrightarrow\quad
 f(x)\in f\left(\bigcup_{i\in I}X_i\right).
\]

\medskip
\noindent $b)$ Dacă
$f(x)\in f\left(\bigcap_{i\in I}X_i\right)$, atunci
$x\in X_i$ pentru orice $i\in I$, de unde
\[
 f(x)\in f(X_i)\quad(\forall i\in I)
 \quad\Longrightarrow\quad
 f(x)\in\bigcap_{i\in I}f(X_i).
\]
\noindent Dacă $f$ este injectivă, din
$f(x)\in\bigcap_{i\in I}f(X_i)$ rezultă
$f(x)\in f(X_i)$ pentru orice $i$, deci
$x\in\bigcap_{i\in I}X_i$.

\medskip
\noindent $c)$
% Demonstratia punctului c) continua pe urmatoarea pagina de manuscris.
.
% Fara preambul si fara intrerupere fortata de pagina.
% NOTE EDITORIALE:
% - Fotografia este numerotata cu 3 in coltul din dreapta jos, desi
%   cererea mentioneaza pagina 2. Acest fragment este deci pagina 3.
% - In manuscris, f(X) si f^{-1}(Y) sunt definite la inceputul exercitiului;
%   anumite semne din dreptul cuantificatorilor sunt suprascrise.
% - Exemplul din dreapta despre f:Z->{0,1} contine intervale/notatii
%   greu lizibile. Sunt redate numai citirile provizorii [-1,0] si [1,2].
%   Manuscrisul noteaza rezultatele ca 1 si 0, nu ca multimi singleton.
% - Sub punctul b) exista etape taiate; au fost omise. In sens strict,
%   si demonstratia de la a) foloseste nejustificat acelasi antecedent x
%   pentru egalitatea imaginilor, desi f nu este presupusa injectiva.
% - Egalitatea la b) pentru o familie vida ar necesita in plus f(A)=B;
%   manuscrisul nu discuta aceasta exceptie.
%   afirmatia f(x) in f(intersectie) nu implica x in toate X_i fara
%   injectivitate: e nevoie de un martor potrivit in intersectie.
% - Pagina se termina cu inceputul punctului c) din demonstratie,
%   fara alte randuri; nu a fost inventata continuarea.

\subsection*{Seminar 2 - Pagina 3}

\noindent\fbox{4}\quad Fie $f:A\longrightarrow B$, $X\subseteq A$, $Y\subseteq B$.
\[
 f(X)=\{f(x)\mid x\in X\}, \qquad
 f^{-1}(Y)=\{a\in A\mid f(a)\in Y\}.
\]
\noindent (Preimaginea lui $Y$; fibra lui $Y$.)

\medskip
\noindent Fie $f:A\longrightarrow B$,
$(X_i)_{i\in I}\subseteq\mathcal{P}(A)$,
$(Y_j)_{j\in J}\subseteq\mathcal{P}(B)$.
\begin{enumerate}
\item[a)]
\[
 f\left(\bigcup_{i\in I}X_i\right)
  =\bigcup_{i\in I}f(X_i).
\]
\item[b)]
\[
 f\left(\bigcap_{i\in I}X_i\right)
  \subseteq\bigcap_{i\in I}f(X_i).
\]
\noindent Dacă $f$ este injectivă, avem egalitate.
\item[c)]
\[
 f^{-1}\left(\bigcup_{j\in J}Y_j\right)
  =\bigcup_{j\in J}f^{-1}(Y_j).
\]
\item[d)]
\[
 f^{-1}\left(\bigcap_{j\in J}Y_j\right)
  =\bigcap_{j\in J}f^{-1}(Y_j).
\]
\end{enumerate}

% Exemplul din marginea dreapta a manuscrisului (citire provizorie).
\noindent\underline{Ex.} $f:\mathbb Z\longrightarrow\{0,1\}$,
\[
 f(x)=\begin{cases}
  0, & x>0,\\
  1, & x\leq 0.
 \end{cases}
 \qquad
 f([-1,0])=1,\quad f([1,2])=0.
\]

\medskip
\noindent\underline{Sol.}

\noindent $a)$ Fie
$ f(x)\in f\left(\bigcup_{i\in I}X_i\right)$.
Atunci $x\in\bigcup_{i\in I}X_i$, deci există $i_1\in I$
astfel încât $x\in X_{i_1}$, de unde
\[
 f(x)\in f(X_{i_1})\subseteq\bigcup_{i\in I}f(X_i).
\]
\noindent Rezultă incluziunea $\subseteq$.

\smallskip
\noindent Reciproc, dacă
$f(x)\in\bigcup_{i\in I}f(X_i)$, există $i_1\in I$
a.î.\ $f(x)\in f(X_{i_1})$. Prin urmare,
\[
 x\in X_{i_1}\subseteq\bigcup_{i\in I}X_i
 \quad\Longrightarrow\quad
 f(x)\in f\left(\bigcup_{i\in I}X_i\right).
\]

\medskip
\noindent $b)$ Dacă
$f(x)\in f\left(\bigcap_{i\in I}X_i\right)$, atunci
$x\in X_i$ pentru orice $i\in I$, de unde
\[
 f(x)\in f(X_i)\quad(\forall i\in I)
 \quad\Longrightarrow\quad
 f(x)\in\bigcap_{i\in I}f(X_i).
\]
\noindent Dacă $f$ este injectivă, din
$f(x)\in\bigcap_{i\in I}f(X_i)$ rezultă
$f(x)\in f(X_i)$ pentru orice $i$, deci
$x\in\bigcap_{i\in I}X_i$.

\medskip
\noindent $c)$
% Demonstratia punctului c) continua pe urmatoarea pagina de manuscris.
.
% Fara preambul si fara intrerupere fortata de pagina.
% NOTE EDITORIALE:
% - Fotografia este numerotata cu 3 in coltul din dreapta jos, desi
%   cererea mentioneaza pagina 2. Acest fragment este deci pagina 3.
% - In manuscris, f(X) si f^{-1}(Y) sunt definite la inceputul exercitiului;
%   anumite semne din dreptul cuantificatorilor sunt suprascrise.
% - Exemplul din dreapta despre f:Z->{0,1} contine intervale/notatii
%   greu lizibile. Sunt redate numai citirile provizorii [-1,0] si [1,2].
%   Manuscrisul noteaza rezultatele ca 1 si 0, nu ca multimi singleton.
% - Sub punctul b) exista etape taiate; au fost omise. In sens strict,
%   si demonstratia de la a) foloseste nejustificat acelasi antecedent x
%   pentru egalitatea imaginilor, desi f nu este presupusa injectiva.
% - Egalitatea la b) pentru o familie vida ar necesita in plus f(A)=B;
%   manuscrisul nu discuta aceasta exceptie.
%   afirmatia f(x) in f(intersectie) nu implica x in toate X_i fara
%   injectivitate: e nevoie de un martor potrivit in intersectie.
% - Pagina se termina cu inceputul punctului c) din demonstratie,
%   fara alte randuri; nu a fost inventata continuarea.

\subsection*{Seminar 2 - Pagina 3}

\noindent\fbox{4}\quad Fie $f:A\longrightarrow B$, $X\subseteq A$, $Y\subseteq B$.
\[
 f(X)=\{f(x)\mid x\in X\}, \qquad
 f^{-1}(Y)=\{a\in A\mid f(a)\in Y\}.
\]
\noindent (Preimaginea lui $Y$; fibra lui $Y$.)

\medskip
\noindent Fie $f:A\longrightarrow B$,
$(X_i)_{i\in I}\subseteq\mathcal{P}(A)$,
$(Y_j)_{j\in J}\subseteq\mathcal{P}(B)$.
\begin{enumerate}
\item[a)]
\[
 f\left(\bigcup_{i\in I}X_i\right)
  =\bigcup_{i\in I}f(X_i).
\]
\item[b)]
\[
 f\left(\bigcap_{i\in I}X_i\right)
  \subseteq\bigcap_{i\in I}f(X_i).
\]
\noindent Dacă $f$ este injectivă, avem egalitate.
\item[c)]
\[
 f^{-1}\left(\bigcup_{j\in J}Y_j\right)
  =\bigcup_{j\in J}f^{-1}(Y_j).
\]
\item[d)]
\[
 f^{-1}\left(\bigcap_{j\in J}Y_j\right)
  =\bigcap_{j\in J}f^{-1}(Y_j).
\]
\end{enumerate}

% Exemplul din marginea dreapta a manuscrisului (citire provizorie).
\noindent\underline{Ex.} $f:\mathbb Z\longrightarrow\{0,1\}$,
\[
 f(x)=\begin{cases}
  0, & x>0,\\
  1, & x\leq 0.
 \end{cases}
 \qquad
 f([-1,0])=1,\quad f([1,2])=0.
\]

\medskip
\noindent\underline{Sol.}

\noindent $a)$ Fie
$ f(x)\in f\left(\bigcup_{i\in I}X_i\right)$.
Atunci $x\in\bigcup_{i\in I}X_i$, deci există $i_1\in I$
astfel încât $x\in X_{i_1}$, de unde
\[
 f(x)\in f(X_{i_1})\subseteq\bigcup_{i\in I}f(X_i).
\]
\noindent Rezultă incluziunea $\subseteq$.

\smallskip
\noindent Reciproc, dacă
$f(x)\in\bigcup_{i\in I}f(X_i)$, există $i_1\in I$
a.î.\ $f(x)\in f(X_{i_1})$. Prin urmare,
\[
 x\in X_{i_1}\subseteq\bigcup_{i\in I}X_i
 \quad\Longrightarrow\quad
 f(x)\in f\left(\bigcup_{i\in I}X_i\right).
\]

\medskip
\noindent $b)$ Dacă
$f(x)\in f\left(\bigcap_{i\in I}X_i\right)$, atunci
$x\in X_i$ pentru orice $i\in I$, de unde
\[
 f(x)\in f(X_i)\quad(\forall i\in I)
 \quad\Longrightarrow\quad
 f(x)\in\bigcap_{i\in I}f(X_i).
\]
\noindent Dacă $f$ este injectivă, din
$f(x)\in\bigcap_{i\in I}f(X_i)$ rezultă
$f(x)\in f(X_i)$ pentru orice $i$, deci
$x\in\bigcap_{i\in I}X_i$.

\medskip
\noindent $c)$
% Demonstratia punctului c) continua pe urmatoarea pagina de manuscris.
